\subsection{Равноускоренное движение}
%%%%%%%%%%%%%%%%%%%%%%%%%%%%%%%%%%%%%%%%%%%%%%%%%%%%
%%%%%%%%%%%%%%%%%%%%%%%%%%%%%%%%%%%%%%%%%%%%%%%%%%%%
\z{В момент, когда опоздавший пассажир вбежал на платформу, мимо него за время $t_1$ прошел предпоследний вагон поезда. Последний вагон прошел мимо пассажира за время $t_2$. На сколько времени пассажир опоздал к отходу поезда? Поезд движется равноускоренно, длина вагонов одинакова.}
%
{$t=\frac{t_2^2+2t_1t_2-t_1^2}{2(t_1-t_2)}$}

%%%%%%%%%%%%%%%%%%%%%%%%%%%%%%%%%%%%%%%%%%%%%%%%%%%%

 \z{Частица, покинув источник, пролетает с постоянной скоростью расстояние $L$, а затем тормозится с ускорением $a$. При какой скорости частицы время движения от ее вылета до остановки будет наименьшим?}
%
{$v=\sqrt{aL}$}
	
%%%%%%%%%%%%%%%%%%%%%%%%%%%%%%%%%%%%%%%%%%%%%%%%%%%%
\z{Тело начинает движение из точки $A$ и движется сначала равноускоренно в течение времени $t_0$, затем с тем же по модулю ускорением — равнозамедленно. Через какое время от начала движения тело вернется в точку $A$?}
%
{$t=(2+\sqrt{2})t_0$}

%%%%%%%%%%%%%%%%%%%%%%%%%%%%%%%%%%%%%%%%%%%%%%%%%%%%	
\z{Тело падает с высоты $h = 45\,\text{м}$. Найдите среднюю скорость $v_\text{ср}$ его движения на второй половине пути. Начальная скорость тела равна нулю. Сопротивление воздуха не учитывать.}
%
{$v_\text{ср}=\sqrt{gh}\frac{\sqrt{2}+1}{2}\approx 25.4\, \nicefrac{\text{м}}{\text{с}}$}

%%%%%%%%%%%%%%%%%%%%%%%%%%%%%%%%%%%%%%%%%%%%%%%%%%%%
\z{Космический корабль движется в открытом космосе со скоростью $V$. Требуется изменить направление скорости на $90\degree$, оставив величину скорости неизменной. Найдите минимальное время, необходимое для такого манёвра, если двигатель может сообщать кораблю в любом направлении ускорение, не превышающее $a$. По какой траектории будет при этом двигаться корабль? }
%
{$t_{min}=\frac{\sqrt{2}v}{a}$;~ по параболе}

